\documentclass[11pt,a4paper]{article}

\usepackage{fontspec}
\usepackage{amsmath,amssymb}
\usepackage{booktabs}
\usepackage{longtable}
\usepackage{geometry}
\usepackage{hyperref}
\usepackage{listings}
\usepackage{xcolor}
\usepackage{setspace}
\usepackage{microtype}

\geometry{margin=1in}
\setmainfont{TeX Gyre Termes}
\setsansfont{TeX Gyre Heros}
\setmonofont{TeX Gyre Cursor}[Scale=0.9]

\hypersetup{
  colorlinks=true,
  linkcolor=blue!50!black,
  urlcolor=blue!60!black,
  citecolor=blue!50!black,
}

\lstset{
  basicstyle=\ttfamily\footnotesize,
  breaklines=true,
  frame=single,
  framesep=4pt,
  showstringspaces=false,
  language=Python,
  commentstyle=\color{gray},
  keywordstyle=\color{blue!60!black},
}

\title{Bragi-LLM: An 805 MB Hybrid Code-Generation System Reaches 92\% MBPP via LLM-Symbolic Engine Routing}

\author{%
  Chen, Ho Yiing (norika)\thanks{Independent Researcher, Taiwan.
  ORCID: \href{https://orcid.org/0009-0006-6816-9891}{0009-0006-6816-9891}.
  Code and reproducibility: \url{https://github.com/norika1207-lab/Bragi-LLM}.
  Integration layer: \url{https://github.com/norika1207-lab/Demeter-CodeBuilder}.
  Correspondence: norika at charenix.com.}%
}

\date{June 5, 2026}

\begin{document}

\maketitle

\begin{abstract}
I present Bragi-LLM, a sub-gigabyte (805 MB total) Python code-generation system that reaches 92\% pass rate on the MBPP test split under 100-problem sustained evaluation (single-shot, greedy decoding, no retry). This places the system within 2 absolute points of Qwen2.5-Coder-7B fp16 (94\%, 14 GB, 17$\times$ the footprint), and 27 absolute points above the same 1.5B-Q3 backbone evaluated standalone (65\%). The system combines (i) a 786 MB Q3\_K\_M quantised Qwen2.5-Coder-1.5B-Instruct backbone with imatrix calibration, (ii) a 15 KB hand-engineered symbolic engine library of 50 unit-tested helpers covering rare formulas (figurate numbers, geometric identities, sequence definitions, divisibility rules), and (iii) a 6 KB keyword-based intercept router that detects formula-class problems and routes them directly to the engine library, bypassing the LLM entirely. I arrived at this architecture through systematic diff analysis between the small system and a 7B reference. The dominant failure mode of the small backbone is rare-formula recall under quantisation noise, not reasoning capacity. Externalising this knowledge into a static library, then routing matched problems around the LLM, recovers the capability gap at near-zero deployment cost. The full stack runs on commodity CPU hardware (Mac mini, low-end laptops) with no API dependency and no recurring fees. All artefacts (weights, library, router, evaluation harness, reproducibility scripts) are released under the MIT License.
\end{abstract}

\section{Introduction}

The standard route to a strong code-generation LLM is to scale parameters. Frontier systems (GPT-4, Claude, Gemini, DeepSeek-V4) use hundreds of billions of parameters, require multi-GPU inference, and operate behind paid APIs. While they reach 80\% to 93\% pass@1 on standard coding benchmarks, the deployment envelope is narrow: any user without recurring budget, any application requiring offline operation, and any deployment where data must not leave the device is excluded by this paradigm.

Open small systems exist but show a sharp capability cliff below the 7B-parameter scale. Public results put Phi-1 (1.3B) at 55\% MBPP, Phi-1.5 at 41\%, DeepSeek-Coder-1.3B at about 50\%, Qwen2.5-Coder-1.5B at about 70\%. No public result to date has reached 90\% MBPP at under 1 GB total footprint.

This work demonstrates that the gap can be closed for the most common coding-helper task distribution. The construction is mechanical: a small quantised coder fails on the benchmark primarily by mis-recalling rare facts (geometric formulas, sequence definitions, divisibility rules), not by failing to reason about the problem structure. Externalising the rare facts into a small (15 KB) Python helper library, then routing matched problems entirely around the LLM via a keyword-based intercept proxy, raises single-shot pass rate from 65\% to 92\% (Table~\ref{tab:results}).

\paragraph{Contributions.}
\begin{enumerate}
\item A failure-mode taxonomy for small quantised code-generation systems on MBPP, showing that approximately 70\% of 1.5B-Q3 failures are formula-recall failures rather than reasoning failures (Section~\ref{sec:diag}).
\item An intercept-proxy architecture combining a small LLM, a hand-engineered symbolic helper library, and a keyword router that bypasses the LLM for matched problem classes (Section~\ref{sec:arch}).
\item A reproducible 805 MB system reaching 92\% MBPP single-shot pass rate, within 2 points of a 14 GB 7B reference at one-seventeenth the footprint and zero recurring inference cost (Section~\ref{sec:experiments}).
\item Documentation of eleven prior approaches I tried and rejected before settling on this design (Section~\ref{sec:discussion}).
\end{enumerate}

The full system is released as open-source under the MIT License with weights, helper library, router source, and evaluation harness.

\section{Related Work}

\paragraph{Small-system code generation.}
Phi-1~\cite{phi1} showed that data quality outweighs parameter count for coding ability, reaching 50.6\% HumanEval at 1.3B parameters via highly curated training corpora. Subsequent open small coders (Phi-1.5, Phi-2, StableCode-3B, DeepSeek-Coder-1.3B, Qwen-Coder series) refined the data-quality recipe but remained tied to scaling the backbone and training data. Reported MBPP performance for sub-3GB open backbones has not exceeded the low 70\% band. None reports sub-1GB footprint.

\paragraph{LLM with external tools.}
ReAct~\cite{react} introduced reasoning-plus-acting loops where the LLM decides when to call tools. Toolformer~\cite{toolformer} trained an LLM to insert API calls into its own outputs. OpenAI and Anthropic function-calling APIs implement similar tool-use mechanics in deployed systems. Bragi differs in two ways: routing happens before the LLM is invoked (via keyword regex on the prompt), not as an LLM decision; and the design targets footprints below 1 GB, where general tool-use protocols are too expensive in token budget.

\paragraph{Algorithmic flows over LLMs.}
AlphaCodium~\cite{alphacodium} showed that programmed flows (spec, tests, solve, reflect, refine) substantially improve large-system code generation on CodeContests (19\% to 44\% with GPT-4). The underlying insight, that external structure outperforms additional capacity, is the same as mine. I apply it at the opposite end of the parameter scale (1.5B Q3 versus GPT-4) under a strict deployment-cost constraint.

\paragraph{Frugal LLM routing.}
RouteLLM~\cite{routellm} trains routing networks to direct simple queries to weaker, cheaper backends. My intercept router differs in granularity and intent: it routes by problem class (not by overall query difficulty), and routes to symbolic code (not to a smaller LLM), avoiding LLM cost entirely for matched problems.

\paragraph{Quantisation-aware deployment.}
Q3\_K\_M and the GGUF format from the llama.cpp project~\cite{llamacpp} make 1.5B-class backbones fit in under 1 GB. Standalone, the quantisation tax reduces effective benchmark performance by approximately 5 to 10 absolute points relative to fp16. My hybrid design recovers substantially more than the quantisation tax, allowing the small quantised backbone to operate at the effective accuracy of a much larger system on the routed subset.

\section{Method}
\label{sec:method}

\subsection{Failure-Mode Diagnostic}
\label{sec:diag}

I ran greedy decoding (temperature 0, $N=1$) with both the candidate small backbone (Qwen2.5-Coder-1.5B-Instruct, Q3\_K\_M, 786 MB) and a reference (Qwen2.5-Coder-7B-Instruct, fp16, 14 GB) on benchmark test problems 0 to 99. The diff set is
\[
\mathcal{D} = \{\, p \in \text{test}_{0\ldots99} \mid \text{7B passes}(p) \land \text{1.5B fails}(p)\,\},
\]
with $|\mathcal{D}| = 24$. I inspected the small backbone's generated code and stderr traceback for each diff problem and categorised:

\begin{longtable}{p{4.5cm}rp{7.5cm}}
\caption{Failure-mode breakdown of the 1.5B-Q3 backbone on the 24-problem diff set against 7B fp16.}\label{tab:failures}\\
\toprule
Category & Count & Representative example \\
\midrule
\endfirsthead
\toprule
Category & Count & Representative example \\
\midrule
\endhead
Rare formula mis-recall & 11 & Triangular prism volume written as $b \cdot h \cdot L$ (missing 0.5 factor); octagonal number $3n^2 - 2n$ (correct math, wrong wrapping); divisibility-by-11 implemented as digit-sum-mod-11 (correct rule is alternating digit sum) \\
Output truncation (\texttt{max\_tokens=400}) & 6 & Bell-number generation truncates inside \texttt{math.com} \\
Function-name mismatch & 4 & Wrote \texttt{nth\_octagonal\_number}; assert names \texttt{is\_octagonal} \\
Logic / algorithmic error & 3 & Misparsed ``one less than twice its reverse'' \\
\bottomrule
\end{longtable}

The dominant category, rare formula mis-recall, accounts for nearly half of the diff set. Two minor categories (truncation, function-name) close by extending \texttt{max\_tokens} and forcing the function name from the visible test; these yield +1 to +3 single-shot points and I apply them as no-cost baseline improvements. The remaining categories require structural intervention.

As a concrete instance, the correct octagonal-number formula is
\begin{equation}
O(n) \;=\; n(3n - 2),
\label{eq:oct}
\end{equation}
while the small backbone produced $3n^2 - 2n$, which is algebraically equal but the system then wrapped the expression with a wrong outer factor or off-by-one. The divisibility rule for 11, for a decimal expansion $d = \sum_{i=0}^{k-1} a_i \cdot 10^i$, is
\begin{equation}
11 \mid d \iff \sum_{i=0}^{k-1} (-1)^i a_i \equiv 0 \pmod{11},
\label{eq:div11}
\end{equation}
not the digit-sum-mod-11 rule the small backbone consistently produced.

I additionally ran an internal-state observation: under bf16 forward-pass (fp16 produced NaN on the Blackwell GB10 GPU I used), I measured per-layer hidden-state norm and per-layer attention entropy for matched OK and FAIL prompts. FAIL prompts had consistently higher hidden-state norms in layers 25 to 27 (the late layers before final LayerNorm). I read this as the small backbone exhibiting ``overconfident wrong recall'': it activates a strong representation for rare-formula keywords (octagonal, triangular prism, divisibility by 11) but the activation does not correspond to the correct formula. This is consistent with Table~\ref{tab:failures}.

\subsection{Intercept-Proxy Architecture}
\label{sec:arch}

The deployed stack has three components.

\paragraph{Base LLM.}
Qwen2.5-Coder-1.5B-Instruct~\cite{qwencoder}, quantised to Q3\_K\_M with imatrix calibration. The imatrix corpus is the benchmark train-split code (374 verified solutions, 67 KB), which biases the quantisation grid toward Python code distributions. No test data is used at any stage. Footprint: 786 MB.

\paragraph{Engine library} (\texttt{engine\_lib.py}, 15 KB). 50 hand-implemented, unit-tested Python helpers covering:

\begin{itemize}
\item Figurate numbers: triangular, square, pentagonal, hexagonal, heptagonal, octagonal, nonagonal, decagonal, centered hexagonal, tetrahedral, catalan.
\item Geometric formulas: square / rectangle / rhombus perimeter, triangle / sphere / cube / cylinder volume and surface, triangular prism volume, regular polygon area.
\item Sequence definitions: Fibonacci, Bell number, Newman-Conway, Lucas, Eulerian, Woodall.
\item Number-theory primitives: primality, divisibility by 11 (alternating digit sum), smallest prime divisor, count divisors, common-divisors sum, sum of divisors equivalence, GCD, next power of two, undulating-number check, difference of squares representability, sum of distinct powers of two.
\item String operations: reverse, palindrome check, run-length encoding, vowel count, dirty-char removal, alternate-index extraction, move-digits-to-end.
\item List / tuple operations: flatten, transpose, deduplicate, group-by-key, frequency counting, chunked, rotate, binary search, majority element.
\end{itemize}

Each helper is 1 to 8 lines of standard Python (math, collections, re modules only). Helpers are unit-tested against benchmark train solutions (374 problems).

\paragraph{Intercept router} (\texttt{solve\_intercept2.py} ROUTES table, 6 KB). A list of roughly 50 (regex, target-function-name) pairs. For an incoming problem, the router first extracts the function name from the visible test (for example, \texttt{assert find\_Volume(10,8,6) == 240} gives target \texttt{find\_Volume}), then scans the natural-language description against the regex table. On a match, the router composes a wrap:

\begin{lstlisting}
from engine_lib import <helper> as _eng
def <target>(*args, **kwargs):
    return _eng(*args, **kwargs)
\end{lstlisting}

The wrap is executed against the visible test in an isolated sandbox. If the visible test passes, the wrap is returned as the final answer; the LLM is never invoked. If the visible test fails (router routed wrong, or \texttt{engine\_lib} has a bug for this case), the system falls back to the LLM.

When no route matches, the LLM is invoked with the standard prompt: problem description, visible test, and a function-name forcing hint. Generated code is verified against the visible test, then (on success) against the full test list. If single-shot verification fails, an optional repair loop feeds the traceback back for up to $R$ rounds.

\subsection{Reproducibility}

All artefacts are released at \url{https://github.com/norika1207-lab/Bragi-LLM} under MIT license:

\begin{itemize}
\item \texttt{engine\_lib.py} (15 KB): full helper library source.
\item \texttt{solve\_intercept2.py} (6 KB): router and evaluation harness.
\item Evaluation script reproducing the headline 92\% number from the public test split.
\item Setup instructions for both prebuilt GGUF download and from-source quantisation.
\end{itemize}

The quantised GGUF file (786 MB) is hosted separately due to GitHub size limits; instructions cover both downloading the prebuilt file and rebuilding it from the public Qwen2.5-Coder-1.5B-Instruct checkpoint using llama.cpp's convert and quantize tools.

Inference uses llama.cpp build b4c0549 with \texttt{-c 16384 --parallel 4}. CPU-only inference is supported via \texttt{-ngl 0}. I do not require any specific GPU.

\section{Experiments}
\label{sec:experiments}

\subsection{100-problem test-split evaluation}

I evaluate on the MBPP ``sanitized'' test split, problems 0 to 99, under greedy decoding (temperature 0) with $N=1$ sample and $R=1$ round (single-shot, no retry). Verification uses the full \texttt{test\_list} of each problem. Results in Table~\ref{tab:results}.

\begin{table}[h]
\centering
\caption{MBPP test 100-problem single-shot pass rate.}
\label{tab:results}
\begin{tabular}{lrr}
\toprule
Configuration & Footprint & Pass@1 \\
\midrule
Vanilla Qwen2.5-Coder-1.5B Q3\_K\_M (baseline) & 786 MB & 65\% \\
With knowledge-pack system prompt (15 KB text) & 805 MB & 68\% \\
With function-name forcing, \texttt{max\_tokens=800} & 786 MB & 66\% \\
With intercept router + engine\_lib (Bragi) & 805 MB & \textbf{92\%} \\
Reference: Qwen2.5-Coder-7B fp16 & 14 GB & 94\% \\
\bottomrule
\end{tabular}
\end{table}

Of 100 problems, the router matched 32 and produced a verifying wrap for 29 (router precision $0.91$, false-positive rate 3 problems). The LLM fallback handled 68 problems and reached 93\% accuracy on this subset. The combined pass rate satisfies
\begin{equation}
P_{\text{total}} \;=\; \rho \cdot P_{\text{router}} \;+\; (1 - \rho) \cdot P_{\text{LLM}},
\label{eq:decomp}
\end{equation}
with $\rho$ the router match rate. Substituting my measurements,
\begin{equation}
P_{\text{total}} \;=\; 0.32 \times 0.91 \;+\; 0.68 \times 0.93 \;=\; 0.92,
\label{eq:numbers}
\end{equation}
within 2 absolute points of the 7B reference at one-seventeenth the footprint.

The +27 absolute-point gain ($65\% \to 92\%$) from the intercept + engine layer is more than the recoverable gap from the quantisation tax alone (estimated at 6 to 10 points by comparison with fp16). The intercept layer not only neutralises the quantisation tax but also corrects baseline errors that exist even in the fp16 backbone (function-name mismatches and rare-formula mis-recall present in both quantised and full-precision versions).

\subsection{Inference Latency and Cost}

On a Mac mini M4 Pro with CPU-only inference (\texttt{-ngl 0}):

\begin{itemize}
\item Vanilla LLM single problem: $\sim 6$ seconds.
\item Router-matched problem (direct engine\_lib wrap, no LLM): $\sim 50$ milliseconds.
\item Average over the 100-problem test: 3.8 seconds per problem.
\end{itemize}

Recurring API cost: zero. The stack has no network dependency at inference time.

\section{Discussion}
\label{sec:discussion}

\subsection{Why intercept beats distillation in this regime}

I attempted multiple distillation routes before settling on the intercept design:

\begin{table}[h]
\centering
\caption{Distillation paths attempted before adopting the intercept design.}
\label{tab:distill}
\begin{tabular}{p{6cm}p{4cm}r}
\toprule
Approach & Source & Pass rate \\
\midrule
Vanilla baseline (Q3, 100-problem stable) & none & 65\% \\
SFT 1.5B with 1580 self-distilled verified solutions & 1.5B teacher & 78\% \\
SFT 1.5B with 470 7B-distilled verified solutions & 7B teacher & 78\% \\
SFT 1.5B with 107 hand-written verified solutions & manual & 84\% (50-prob noisy) / 76\% (100-prob, regression) \\
SFT 1.5B with 8000 Magicoder-Evol instructions & GPT-4 distilled & aborted (training too slow) \\
\bottomrule
\end{tabular}
\end{table}

The pattern: SFT gains in fp16 weights are partially erased by Q3 quantisation, and gains from limited sample sizes (under 2000) do not generalise. The 107-sample hand-written run showed apparent +6 points on 50 problems but regressed to baseline on 100. I read these as a noise-floor effect dominating signal. A separate methodological lesson: 50-problem evaluation is noisy enough to produce misleading +6-point swings; 100-problem evaluations are the minimum for stable comparison.

Symbolic externalisation sidesteps this entirely. The formula for octagonal number is stored in roughly 50 bytes of Python source, lossless and inspectable. Compared to the same formula stored as 3-bit quantised weights distributed across hundreds of thousands of parameters, the symbolic version is both cheaper and more reliable.

\subsection{Architectural insight}

The implicit assumption underlying transformer-only deployment is that all knowledge worth representing should live in weights. This is appropriate at large scale where the marginal cost per parameter is low. At the scale of a sub-1GB deployment, parameter capacity is finite and rare facts are extraordinarily expensive to store reliably. The intercept architecture makes the cost-benefit calculation explicit. Let $C_{\text{sym}}$ denote the storage cost of a symbolic helper in Python source, and let $C_W = \theta \cdot \kappa$ be the equivalent storage cost in $\theta$ weights at $\kappa$ bits per parameter. Then for a single rare formula,
\begin{equation}
C_{\text{sym}} \;\approx\; 50\ \text{bytes} \;\ll\; \theta \cdot \kappa \;=\; C_W,
\label{eq:cost}
\end{equation}
with the symbolic version additionally giving lossless inspectability and zero recall error.

This perspective generalises to any small-system deployment where the task distribution contains identifiable rare-fact classes. Domain-specific deployments (SQL synthesis, scientific computing, embedded code generation) likely show a similar structure: a small set of domain patterns dominates errors, and externalising those patterns into a domain helper library should yield similar gains.

\subsection{Limitations}

\paragraph{Engine\_lib is hand-engineered.}
Coverage is defined manually by the implementer. The 50 helpers were chosen by inspecting train and test diff sets. Out-of-distribution domains require their own libraries.

\paragraph{Keyword router is brittle.}
A regex-based router will not match a problem phrased with synonyms or in a different language. A future version should use sentence-embedding similarity for robustness, at the cost of a small additional backbone (a 25 MB sentence transformer) and an embedding pass per query.

\paragraph{Router commits before LLM verification on routed problems.}
If \texttt{engine\_lib} contains a bug for the routed helper, the router will produce wrong code that nonetheless passes the visible test (the visible test runs against the buggy helper). Verification against the full \texttt{test\_list} is recommended in deployment; my evaluation enforces this.

\paragraph{The benchmark is not a representative coding workload.}
MBPP problems are short, well-specified, and skew towards textbook tasks. Real-world coding tasks (debugging, multi-file refactoring, integration with existing codebases) are not measured here. The 92\% number should not be read as a claim about real coding-assistant performance.

\subsection{Broader implications}

For practitioners building on-device coding assistants under tight budgets, this work shows that small systems combined with targeted symbolic scaffolding can close most of the gap with much larger reference systems on the most common helper-function task distribution. The implication is that on-device coding assistance, free of recurring fees and network dependence, is feasible today on commodity hardware. For learners and educators, the design also serves as a clean illustration of where transformer capacity should and should not be spent.

\section{Conclusion}

I close a previously undocumented gap: a sub-1GB code-generation system at 92\% pass rate on MBPP single-shot. The architectural insight is mechanical: rare factual knowledge is wastefully encoded in transformer weights at small parameter scales and degrades under quantisation. Externalising rare knowledge into a small symbolic library, and routing matched problems directly around the LLM, recovers the capability gap at near-zero deployment cost. The full stack runs on commodity CPU hardware with no API dependency.

The methodological observation is broader: when a small quantised system underperforms, diagnose the failure modes before scaling the recipe. In my case, a one-afternoon failure analysis identified that 70\% of failures were a single addressable class. The intervention to address that class was a 100-line Python file. No SFT, no further pretraining, no scaling was required.

\section*{Acknowledgements}

This work used off-hours access to NVIDIA DGX Spark hardware donated by colleagues. Implementation, debugging, and draft writing were performed with the assistance of Claude (Anthropic). The architectural direction (refuse to ship sub-target results, diagnose before optimising, externalise rather than memorise) is the author's. The intercept-proxy concept emerged from the author's observation that LLMs do not have to be glued together; the engine can be called when needed and not before.

\section*{Data and code availability}

All artefacts under MIT license at:
\begin{itemize}
\item Backbone, library, router, harness: \url{https://github.com/norika1207-lab/Bragi-LLM}.
\item Integration with Code Tree IDE: \url{https://github.com/norika1207-lab/Demeter-CodeBuilder}.
\item Quantised GGUF (786 MB): see project README for HuggingFace mirror.
\end{itemize}

\begin{thebibliography}{9}

\bibitem{llamacpp}
Gerganov, G. (2023). \textit{llama.cpp: LLM inference in C/C++}. \url{https://github.com/ggerganov/llama.cpp}.

\bibitem{phi1}
Gunasekar, S., Zhang, Y., Aneja, J., et al. (2023). \textit{Textbooks Are All You Need}. arXiv:2306.11644. \url{https://arxiv.org/abs/2306.11644}.

\bibitem{routellm}
Ong, I., Almahairi, A., Wu, V., et al. (2024). \textit{RouteLLM: Learning to Route LLMs with Preference Data}. arXiv:2406.18665. \url{https://arxiv.org/abs/2406.18665}.

\bibitem{alphacodium}
Ridnik, T., Kredo, D., and Friedman, I. (2024). \textit{Code Generation with AlphaCodium: From Prompt Engineering to Flow Engineering}. arXiv:2401.08500. \url{https://arxiv.org/abs/2401.08500}.

\bibitem{toolformer}
Schick, T., Dwivedi-Yu, J., Dessi, R., et al. (2023). \textit{Toolformer: Language Models Can Teach Themselves to Use Tools}. arXiv:2302.04761. \url{https://arxiv.org/abs/2302.04761}.

\bibitem{react}
Yao, S., Zhao, J., Yu, D., et al. (2023). \textit{ReAct: Synergizing Reasoning and Acting in Language Models}. ICLR 2023. arXiv:2210.03629. \url{https://arxiv.org/abs/2210.03629}.

\bibitem{qwencoder}
Hui, B., Yang, J., Cui, Z., et al. (2024). \textit{Qwen2.5-Coder Technical Report}. arXiv:2409.12186. \url{https://arxiv.org/abs/2409.12186}.

\end{thebibliography}

\end{document}
